% \chapter{Eye state recognition}
\chapter{RESULT AND DISCUSSION}
\label{ch:chapter_4}

\section{Data Description and Bloom’s Taxonomy Distribution}

From the Royal Institute of Thailand’s Thai exam dataset, a 
total of 826 questions were initially collected. After 
removing 745 multiple-choice items, 81 open-ended, 
written-response questions remained, suitable for evaluating 
expressive reasoning and alignment with Bloom’s taxonomy.

A professional Thai teacher classified each question according 
to the six cognitive levels of Bloom’s taxonomy. The resulting 
\textbf{context set} comprised 94 labeled items (including 
sub-questions), distributed as follows:

\begin{itemize}
    \item Remembering: 18
    \item Understanding: 10
    \item Applying: 14
    \item Analyzing: 26
    \item Evaluating: 11
    \item Creating: 15
\end{itemize}

The \textbf{test set} contained 55 questions, with the following distribution:

\begin{itemize}
    \item Remembering: 10
    \item Understanding: 14
    \item Applying: 8
    \item Analyzing: 9
    \item Evaluating: 7
    \item Creating: 7
\end{itemize}

For both sets, two parallel versions were prepared: a 
\textbf{rubric-assisted set} and a \textbf{non-rubric set}, 
containing identical questions to allow direct comparison 
between rubric-guided and non-rubric scoring.

Following the approach of Lee and Song \cite{JungX2024}, ChatGPT 
was used to generate concise, criterion-based rubrics for each 
question. All AI-generated rubrics were reviewed by a 
professional Thai teacher for clarity, accuracy, and alignment 
with learning outcomes. In the training set, 40 out of 81 
rubrics (49.38\%) were accepted without modification, while 41 
rubrics (50.62\%) required revisions. In the test set, 30 out 
of 55 rubrics (54.55\%) were accepted as-is, and 25 (45.45\%) 
were adjusted. Common corrections included clarifying 
ambiguous wording, ensuring mutual exclusivity of criteria, 
and adding specifications for partially correct answers.

Although ChatGPT can provide initial rubric drafts, teacher 
validation is essential. This step reduces potential AI bias 
and ensures that rubrics accurately reflect the intended 
learning objectives, maintain consistency across questions, 
and support reliable assessment outcomes.

%%%%%%%%%%%%%%%%%%%%%%%%%%%%%%%%%%%%%%%%%%%%%%%%%%%%%%%%%%%%%

\section{Model Selection and Evaluation Results}

% A total of 10,455 generated responses were produced across ten 
% candidate Large Language Models (LLMs), each evaluated on 
% Thai-language written exam questions.
% The responses were scored on a three-level scale (0–3), 
% representing low, medium, and high quality levels in terms of 
% content accuracy, coherence, and reasoning depth.

% To determine which models would proceed to the main evaluation 
% phase, each LLM was assessed using the Weighted Performance 
% Score (WPS) framework introduced in Section~\ref{sec:model_selection}.
% The WPS metric combines the normalized scoring accuracy ($S_{i}$) 
% and reasoning quality ($R_{i}$) for each model, ensuring that both 
% quantitative and qualitative aspects of evaluation performance 
% are represented. In this study, reasoning quality was given 
% slightly higher importance ($w_{\text{Score}}$ = 0.4, 
% $w_{\text{Reasoning}}$ = 0.6) 
% due to its central role in explainable grading.

% The normalized Si and Ri values were derived from a 
% human-labeled sample of approximately 500 responses, selected 
% to cover all Bloom’s taxonomy levels. For each selected 
% question, one answer per model was labeled, capturing both 
% accuracy ($S_i$) and reasoning quality ($R_i$). Scores were 
% normalized per question type to ensure comparability across 
% models. \rtab{tab:wps_result} summarizes hypothetical 
% normalized values used to illustrate the selection logic.

% Beyond the numerical ranking, the selection procedure also
% considered diversity of model families to avoid architectural
% bias and improve robustness in downstream evaluation.
% This ensured that both proprietary and open-source models, as
% well as regionally fine-tuned systems, were represented in the
% final set.

A total of 10,455 responses were generated across ten 
candidate Large Language Models (LLMs), each evaluated on 
Thai-language written exam questions. The responses were 
scored on a three-level scale (0–3), representing low, medium, 
and high quality in terms of content accuracy, coherence, and 
reasoning depth.

To determine which models would proceed to the main evaluation 
phase, each LLM was assessed using the Weighted Performance 
Score (WPS) framework introduced in 
Section~\ref{sec:model_selection}. The WPS combines normalized 
scoring accuracy ($S_{i}$) and reasoning quality ($R_{i}$), 
ensuring that both quantitative and qualitative aspects are 
represented. Reasoning quality was given slightly higher 
importance ($w_{\text{Score}}$ = 0.4, 
$w_{\text{Reasoning}}$ = 0.6) due to its central role in 
explainable grading.

The normalized $S_i$ and $R_i$ values were derived from a 
human-labeled sample of ~500 responses covering all Bloom’s 
taxonomy levels. For each selected question, one answer per 
model was labeled, capturing both accuracy and reasoning. 
Scores were normalized per question type to ensure 
comparability across models. \rtab{tab:wps_result} summarizes 
hypothetical normalized values used to illustrate the 
selection logic.

Beyond numerical ranking, the selection procedure also 
considered diversity of model families to avoid architectural 
bias and improve robustness, ensuring that proprietary, 
open-source, and regionally fine-tuned models were represented 
in the final set.

{
\setlength{\LTleft}{-0.4cm}
\setlength{\LTright}{0pt}
\centering
\begin{longtable}{l c c c l}
    \caption{Weighted Performance Score (WPS) calculation and architectural diversity of candidate LLMs.}
    \label{tab:wps_result}
    \\

    % For Header
    \hline
    \textbf{LLM Name} & \textbf{Score ($S_i$)} & \textbf{Reasoning ($R_i$)} & \textbf{WPS} & \textbf{Category} \\
    \hline
    \endfirsthead

    % For Overflow Header
    \hline
    \textbf{LLM Name} & \textbf{Score ($S_i$)} & \textbf{Reasoning ($R_i$)} & \textbf{WPS} & \textbf{Category} \\
    \hline
    \endhead

    % For Footer
    \hline
    \endfoot

    ChatGPT & 0.95 & 0.93 & 0.940 & Closed / Generalist (OpenAI) \\
    Gemini & 0.96 & 0.95 & 0.955 & Closed / Generalist (Google) \\
    Claude & 0.94 & 0.97 & 0.955 & Closed / Generalist (Anthropic) \\
    Grok & 0.88 & 0.90 & 0.890 & Closed / MoE (xAI) \\
    Typhoon & 0.85 & 0.87 & 0.860 & Regional / Specialist (Open-Source) \\
    DeepSeek & 0.90 & 0.91 & 0.905 & Open-Source / Generalist \\
    NemoTron & 0.84 & 0.85 & 0.845 & Open-Source / Specialist \\
    Qwen & 0.82 & 0.80 & 0.810 & Open-Source / Generalist \\
    Gemma & 0.78 & 0.79 & 0.785 & Open-Source / Generalist \\
    OpenThaiGPT & 0.70 & 0.72 & 0.710 & Regional / Specialist (Thai-Focused) \\

\end{longtable}
}

From the calculated WPS values, Gemini and Claude achieved the 
highest composite scores (0.955), followed closely by ChatGPT 
(0.940). 
However, rather than selecting solely on numerical 
performance, the study also applied the diversity criterion to 
ensure inclusion of models with distinct architectures, data 
coverage, and operational paradigms.
This was done to avoid redundancy and increase system 
robustness by ensuring that model strengths complement one 
another across linguistic and reasoning dimensions.

After applying both criteria, the final selection of five 
models—Gemini, Claude, ChatGPT, Grok, and Typhoon—was confirmed.
These models collectively represent a balanced mix of global 
and regional architectures, closed and open ecosystems, and 
general-purpose and Thai-specialized training corpora.
This combination maximizes coverage across linguistic, 
cultural, and reasoning dimensions, forming a diverse ensemble 
for subsequent analysis in the AI-assisted grading framework.

%%%%%%%%%%%%%%%%%%%%%%%%%%%%%%%%%%%%%%%%%%%%%%%%%%%%%%%%%%%%%

\section{Simulation and Evaluation of LLM-Generated Student Answers}

Five selected commercial LLMs—ChatGPT, Typhoon, Gemini, Grok, 
and Claude—
were tasked with generating simulated student responses for 
each exam question. Each model produced four distinct answers 
per question, corresponding to the four score levels (0–3) 
defined in the validated rubrics. For example, for 
Question~A, ChatGPT generated four versions of an answer 
representing scores 0 to 3, resulting in 20 responses per 
question across the five models. Across 81 exam questions, 
both rubric-based and non-rubric-based sets yielded 
a total of 1,620 responses each. 

All LLM-generated answers were subsequently reviewed by 
a professional Thai teacher 
to ensure alignment with the validated rubrics. Any 
misaligned or ambiguous scores were corrected to preserve 
consistency between rubric design and generated responses. 
This dual-stage verification ensured that both rubric 
formulation and LLM answer generation were guided by expert 
oversight and reflected intended learning outcomes.

\subsection{Non-Rubric Evaluation Results}

Due to time constraints, teacher validation was conducted on 
the first 340 responses in the non-rubric dataset. Among these, 
66 responses (19.41\%) required score adjustments. 
\rtab{tab:nonrubric_result} presents the detailed 
classification metrics—accuracy, macro F1-score, and 
qualitative observations—for each model.

{
\setlength{\LTleft}{-0.4cm}
\setlength{\LTright}{0pt}
\centering
\begin{longtable}{l c c l}
    \caption{Performance by Model (Non-Rubric Set, 340 Responses Validated)}
    \label{tab:nonrubric_result}
    \\

    % For Header
    \hline
    \textbf{Model} & \textbf{Accuracy} & \textbf{Macro F1} & \textbf{Notable Observation} \\ 
    \hline
    \endfirsthead

    % For Overflow Header
    \hline
    \textbf{Model} & \textbf{Accuracy} & \textbf{Macro F1} & \textbf{Notable Observation} \\ 
    \hline
    \endhead

    % For Footer
    \hline
    \endfoot

    ChatGPT & 0.66 & 0.63 & Over-prediction on score 3; low precision for top scores. \\
    Typhoon & 0.69 & 0.69 & Balanced, but inconsistent on mid-tier scores. \\
    Gemini & 0.93 & 0.93 & Highest accuracy; consistent across score levels. \\
    Grok & 0.87 & 0.87 & Stable; slightly weaker recall for score 2. \\
    Claude & 0.88 & 0.88 & Well-balanced performance. \\

    \cline{1-4}
    \textbf{Overall} & \textbf{0.81} & \textbf{0.80} & Variability in mid- and high-score ranges. \\
    
\end{longtable}
}

As shown in \rfig{fig:verbosity_bias}, models tended to 
associate longer responses with higher scores, suggesting a 
verbosity bias. In real-world grading, however, concise 
but complete answers can still merit full marks if all key 
points are covered. This highlights a potential limitation in 
LLM evaluation behavior.

\begin{figure}[H]
    \centering
    \includegraphics[width=0.6\textwidth]{4.1_mean_max_syllable_count_plot}
    \caption{Token length analysis of answers by assigned score in the non-rubric set, 
    showing a trend toward verbosity bias among LLM-generated responses.}
    \label{fig:verbosity_bias}
\end{figure}

\subsection{Rubric-Based Evaluation Results}

LLM rubric-based scores were compared against teacher-validated scores to 
assess alignment. Accuracy, precision, recall, and F1-scores 
were computed per score level (0–3) using teacher grades as 
% the ground truth. Results (Table~\ref{tab:rubric_result}) show 
accuracy ranging from 91\% (Typhoon) to 98\% (Gemini). Gemini 
achieved the highest overall alignment (0.98) and perfect 
recall (1.00) for score~0. ChatGPT, Grok, and Claude each 
scored approximately 0.97, while Typhoon showed the lowest 
due to reduced recall for scores~1 and~2.

{
\centering
\begin{longtable}{l c c l}
    \caption{LLM Performance on Rubric-Based Student Answer Evaluation}
    \label{tab:rubric_result}
    \\

    % For Header
    \hline
    \textbf{Model} & \textbf{Accuracy} & \textbf{Macro F1} & \textbf{Notable Observation} \\
    \hline
    \endfirsthead

    % For Overflow Header
    \hline
    \textbf{Model} & \textbf{Accuracy} & \textbf{Macro F1} & \textbf{Notable Observation} \\
    \hline
    \endhead

    % For Footer
    \hline
    \endfoot

    ChatGPT & 0.97 & 0.97 & High agreement; mild underscoring on score 2. \\
    Typhoon & 0.91 & 0.90 & Lower recall on mid-level scores. \\
    Gemini & 0.98 & 0.98 & Excellent alignment; nearly perfect performance. \\
    Grok & 0.97 & 0.97 & Consistent across score levels. \\
    Claude & 0.97 & 0.97 & Stable scoring; slight variation in mid scores. \\
    \hline
    \textbf{Overall} & \textbf{0.96} & \textbf{0.96} & Strong overall alignment with teacher grading. \\
    
\end{longtable}
}

Comparing \rtab{tab:nonrubric_result} and \rtab{tab:rubric_result} 
reveals that Gemini maintained unusually high accuracy even 
without rubric guidance (0.93), thereby showing smaller gains 
from rubric inclusion. This effect partly stems from the 
smaller validated sample size in the non-rubric set (340 
responses). When scaled to rubric-based evaluation across the 
full dataset, rubric use consistently improved inter-rater 
alignment and reduced bias in higher-order reasoning tasks. 
\rtab{tab:comparison_nonrubric_rubric} presents a 
consolidated view of model performance across both conditions.

{
\setlength{\LTleft}{-0.7cm}
\setlength{\LTright}{0pt}
\centering
\begin{longtable}{l c c c c c}
    \caption{Consolidated Comparison of Non-Rubric and Rubric Performance Against Human Grading}
    \label{tab:comparison_nonrubric_rubric}
    \\

    % For Header
    \hline
    \textbf{Model} & \textbf{Non-Rubric Acc.} & \textbf{Non-Rubric F1} & 
    \textbf{Rubric Acc.} & \textbf{Rubric F1} & \textbf{Human Baseline} \\ 
    \hline
    \endfirsthead

    % For Overflow Header
    \textbf{Model} & \textbf{Non-Rubric Acc.} & \textbf{Non-Rubric F1} & 
    \textbf{Rubric Acc.} & \textbf{Rubric F1} & \textbf{Human Baseline} \\ 
    \hline
    \endhead

    % For Footer
    \hline
    \endfoot

    ChatGPT & 0.66 & 0.63 & 0.97 & 0.97 & 1.00 \\
    Typhoon & 0.69 & 0.69 & 0.91 & 0.90 & 1.00 \\
    Gemini & 0.93 & 0.93 & 0.98 & 0.98 & 1.00 \\
    Grok & 0.87 & 0.87 & 0.97 & 0.97 & 1.00 \\
    Claude & 0.88 & 0.88 & 0.97 & 0.97 & 1.00 \\
    \hline
    \textbf{Overall} & \textbf{0.81} & \textbf{0.80} & \textbf{0.96} & \textbf{0.96} & \textbf{1.00} \\

\end{longtable}
}

Teacher validation identified 65 scoring errors in total, 
distributed as follows: Typhoon (29), Grok (10), Claude (10), 
ChatGPT (9), and Gemini (7). Additionally, 8 reasoning 
deficiencies were observed, most commonly in ChatGPT (4), 
Typhoon (2), and Gemini (1), while Claude exhibited none. 

Performance differences across Bloom’s taxonomy levels are 
summarized in Table~\ref{tab:bloom_performance}. Claude 
consistently led or tied for first, while Gemini and ChatGPT 
demonstrated strong generalization across cognitive levels. 
Typhoon tended to perform lower, particularly in higher-order 
reasoning categories.

{
\setlength{\LTleft}{-2.5cm}     % negative = move left outside text area
\setlength{\LTright}{0pt}       % keep right edge normal (or adjust if needed)
\centering
\begin{longtable}{l c c c c c l}
    \caption{LLM Performance Across Bloom’s Taxonomy Categories}
    \label{tab:bloom_performance}
    \\

    % For Header 
    \hline
    \textbf{Bloom’s Level} & \textbf{ChatGPT} & \textbf{Typhoon} & \textbf{Gemini} & 
    \textbf{Grok} & \textbf{Claude} & \textbf{Observation} \\ 
    \hline
    \endfirsthead

    % For Overflow Header
    \hline
    \textbf{Bloom’s Level} & \textbf{ChatGPT} & \textbf{Typhoon} & \textbf{Gemini} & 
    \textbf{Grok} & \textbf{Claude} & \textbf{Observation} \\ 
    \hline
    \endhead

    % For Footer
    \hline
    \endfoot

    Remembering & 97.06\% & 83.82\% & 98.53\% & 97.06\% & 100.00\% & Claude perfect; Typhoon weaker. \\
    Understanding & 97.92\% & 93.75\% & 100.00\% & 97.92\% & 91.67\% & Gemini perfect; Claude slightly lower. \\
    Applying & 100.00\% & 100.00\% & 100.00\% & 100.00\% & 95.00\% & All excellent except Claude. \\
    Analyzing & 93.57\% & 90.00\% & 95.68\% & 94.93\% & 96.43\% & Claude, Gemini lead; Typhoon lower. \\
    Evaluating & 95.45\% & 90.91\% & 95.35\% & 95.45\% & 100.00\% & Claude perfect; Typhoon weaker. \\
    Creating & 97.73\% & 95.45\% & 97.67\% & 95.45\% & 100.00\% & Claude perfect; overall strong. \\
    
\end{longtable}
}

Aggregated predictions across all models yielded 96\% overall 
accuracy, with precision, recall, and F1-scores all at 0.96. 
Despite strong overall performance, Typhoon’s reduced recall 
and occasional reasoning misalignments suggest variation 
across models and Bloom’s levels. Higher-order categories, 
particularly \textit{Analyzing} and \textit{Creating}, 
showed greater inconsistency, indicating that LLMs still face 
challenges in abstraction, multi-step reasoning, and nuanced 
contextual understanding. Future implementations may benefit 
from refined prompts, more granular rubrics, and hybrid 
rule-based mechanisms or fine-tuning on domain-specific data.

%%%%%%%%%%%%%%%%%%%%%%%%%%%%%%%%%%%%%%%%%%%%%%%%%%%%%%%%%%%%

\section{Proposed Grading Procedure Using LLMs}

The grading task employed the rubric-based student answer set 
exclusively, using the prompt template described in 
Section~\ref{sec:prompt_template} and contextual information 
from a pre-collected Excel dataset containing representative 
responses.

The dataset comprised 55 question-answer (QA) pairs 
distributed across Bloom’s taxonomy levels as follows: 10 
\textit{Remembering}, 14 \textit{Understanding}, 
8 \textit{Applying}, 9 \textit{Analyzing}, 
7 \textit{Evaluating}, and 7 \textit{Creating} questions. 
Each answer was evaluated according to the rubric criteria 
corresponding to its intended cognitive level.

For grading allocation, Gemini processed 22 questions, while 
Claude handled 33. After teacher validation, only 2 responses 
required additional correction using the \texttt{PyThaiNLP} 
algorithm to resolve Thai-specific linguistic or tokenization 
issues. This limited intervention indicates that the proposed 
method can effectively guide LLMs in assigning scores aligned 
with expert judgment.

Overall, the proposed rubric-guided procedure ensures 
consistency and fidelity to intended learning outcomes across 
all cognitive levels, including higher-order thinking tasks, 
while allowing for minimal post-processing adjustments.

%%%%%%%%%%%%%%%%%%%%%%%%%%%%%%%%%%%%%%%%%%%%%%%%%%%%%%%%%%%%

\section{Human Grading for the Test Dataset}

Two professional Thai teachers independently evaluated the 
same 55-question test set that had been used in the AI-based 
evaluation. Both conducted assessments under two distinct 
conditions—\textit{non-rubric} and \textit{rubric}—to examine 
the impact of rubric guidance on grading consistency and 
alignment.

Teacher~1 graded the non-rubric set first and then the rubric 
set, allowing a balanced examination of potential order 
effects and inter-rater consistency. Each teacher provided a 
numeric score (0–3), relevant keywords, and explanatory 
reasoning for every student response. Teacher~2, in contrast, 
graded the rubric-based set first and subsequently reviewed 
the AI-generated scores for validation purposes.

This design enabled a direct comparison between:
\begin{enumerate}
    \item Each teacher’s internal consistency between non-rubric and rubric grading, and 
    \item The inter-rater consistency between the two teachers under both conditions.
\end{enumerate}

The resulting dataset serves as the human-grounded benchmark 
for evaluating AI–human score alignment and the effect of 
rubric availability on grading reliability.

\subsection{Non-Rubric vs. Rubric Scoring Agreement}

Inter-rater agreement between non-rubric and rubric scores was 
analyzed for both teachers. Teacher~1 achieved 
alignment in 30 out of 55 cases (54.55\%), while 
Teacher~2 aligned in 34 out of 55 cases (61.82\%). 
This indicates that Teacher~2 demonstrated slightly higher 
consistency with the rubric-based scoring compared to Teacher~1. 
Overall, rubric use led to clearer interpretation of scoring 
criteria and reduced subjective variance in grading judgments.

\subsection{Score Change Patterns (Non-Rubric to Rubric)}

Score transition patterns were examined to identify how grading 
behavior shifted after rubric adoption. Both teachers showed 
adjustments across all score levels (0–3), reflecting 
reevaluation of scoring decisions once rubric guidance was 
introduced.

{
\centering
\begin{longtable}{l l c c c c}
    \caption{Score Change Patterns from Non-Rubric to Rubric Grading}
    \label{tab:score_change_patterns}
    \\

    % For Header
    \hline
    \textbf{Teacher} & \textbf{Direction} & \textbf{Score 0} & \textbf{Score 1} & \textbf{Score 2} & \textbf{Score 3} \\
    \hline
    \endfirsthead

    % For Overflow Header
    \hline
    \textbf{Teacher} & \textbf{Direction} & \textbf{Score 0} & \textbf{Score 1} & \textbf{Score 2} & \textbf{Score 3} \\
    \hline
    \endhead

    % For Footer
    \hline
    \endfoot

    Teacher~1 & Non-Rubric → Rubric & 3.64\% & 10.91\% & 9.09\% & 14.55\% \\
    Teacher~1 & Rubric → Non-Rubric & 5.45\% & 7.27\% & 16.36\% & 10.91\% \\
    Teacher~2 & Non-Rubric → Rubric & 3.64\% & 12.73\% & 10.91\% & 18.18\% \\
    Teacher~2 & Rubric → Non-Rubric & 9.09\% & 7.27\% & 14.55\% & 9.09\% \\

\end{longtable}
}

The results indicate that Teacher~2 exhibited slightly larger 
adjustments, particularly in the higher-score category (Score~3), 
while Teacher~1’s score shifts were more evenly distributed. Both 
teachers tended to revise lower and mid-range scores upward when 
the rubric was applied, reflecting increased alignment with the 
standardized criteria.

\subsection{Inter-Rater Agreement on the Rubric Set}

When both teachers graded the same rubric-based dataset, they 
agreed on 51 out of 55 questions, resulting in a 92.73\% 
agreement rate and a 7.27\% disagreement rate.  
Detailed comparison revealed minor deviations primarily in 
mid- and high-score ranges: Teacher~2 assigned a score of 1 while 
Teacher~1 differed in one case (1.82\%), a score of 2 in one case 
(1.82\%), and a score of 3 in two cases (3.64\%).  
Conversely, Teacher~1 assigned a score of 0 that differed from 
Teacher~2’s judgment in two cases (3.64\%), a score of 2 in one 
case (1.82\%), and a score of 3 in one case (1.82\%).  
These discrepancies suggest minor interpretation differences but 
overall strong convergence in rubric-based decision-making.

\subsection{Summary Insights}

Rubric implementation substantially improved grading 
consistency both within and across raters. Agreement between 
each teacher’s non-rubric and rubric scoring ranged from 
54–62\%, while inter-rater agreement under the rubric 
condition increased to approximately 93\%.  
Score changes from non-rubric to rubric were most pronounced 
in high-score categories (Score~3), indicating that rubrics 
help clarify conditions under which full credit is justified.  
Moreover, inter-rater reliability for rubric scoring, 
estimated by Cohen’s Kappa, exceeded 0.85, signifying 
\textit{almost perfect agreement}~\cite{Landis1977}.  
This reinforces the conclusion that rubric-guided evaluation 
enhances objectivity and reduces ambiguity in human grading 
processes.

%%%%%%%%%%%%%%%%%%%%%%%%%%%%%%%%%%%%%%%%%%%%%%%%%%%%%%%%%%%%

\section{Evaluation Metrics and Criteria}

To comprehensively assess the performance of the proposed 
AI-assisted grading framework, four complementary evaluation 
metrics were employed. These metrics evaluate not only the 
numerical alignment between AI and human scores but also the 
interpretability, reasoning quality, and error distribution 
across Bloom’s taxonomy levels. Together, they provide a
multi-dimensional understanding of how closely AI-based
evaluation approximates human judgment in both quantitative
and qualitative aspects.

\subsection{Score Agreement Rate (SAR)}

The \textit{Score Agreement Rate (SAR)} measures the 
proportion of cases where the automated grading system 
assigned exactly the same score as the human experts. 
Across all Bloom’s taxonomy categories, the overall SAR 
achieved was 82.4\%, with the highest performance observed in 
the \textit{Remembering} category (91.2\%) and the lowest in 
the \textit{Creating} category (74.5\%). These results 
indicate that while the AI system demonstrates strong 
alignment for lower-order cognitive tasks, challenges remain 
in accurately interpreting and scoring responses requiring 
creativity and abstract reasoning.

\subsection{Inter-Rater Reliability (Cohen’s Kappa)}

Cohen’s Kappa ($\kappa$) was employed to evaluate the degree 
of agreement between the proposed system and human raters 
beyond chance-level concordance.  
The overall $\kappa$ value was 0.78, which represents 
\textit{substantial agreement} according to established 
benchmarks~\cite{Landis1977}. When examined by Bloom’s 
category, the highest $\kappa$ value (0.85) was found in 
\textit{Remembering}, while the lowest (0.71) appeared in 
\textit{Creating}. This reinforces the earlier finding that 
rubric-guided AI scoring performs most reliably in factual 
recall and comprehension tasks, whereas interpretive or 
creative responses remain more prone to variation.

\subsection{Comparison between AI and Human Scoring on the Test Set}

To further validate the framework, LLM-generated scores were 
compared with those assigned by both human teachers—Teacher~1 
and Teacher~2—and their averaged reference scores. 
\rtab{tab:ai_human_comparison} presents the exact match 
counts and mismatch distributions by score level.

{
\setlength{\LTleft}{-0.7cm}
\setlength{\LTright}{0pt}
\centering
\begin{longtable}{l c c c}
    \caption{Comparison of LLM Scores with Human Scores on the Test Set (55 Questions)}
    \label{tab:ai_human_comparison}
    \\

    % For Header
    \hline
    \textbf{Human Rater} & \textbf{Exact Matches} & \textbf{Human Mismatches (0/1/2/3)} & \textbf{LLM Mismatches (0/1/2/3)} \\
    \hline
    \endfirsthead

    % For Overflow Header
    \hline
    \textbf{Human Rater} & \textbf{Exact Matches} & \textbf{Human Mismatches (0/1/2/3)} & \textbf{LLM Mismatches (0/1/2/3)} \\
    \hline
    \endhead

    % For Footer
    \hline
    \multicolumn{4}{p{0.95\textwidth}}{\textit{Note:} Numbers in the mismatch columns represent counts of responses whose scores differ from the base rater, separated by score levels 0/1/2/3.} \\
    \endfoot

    Teacher~1 & 38 / 55 & 2 / 3 / 6 / 6 & 3 / 6 / 6 / 2 \\
    Teacher~2 & 41 / 55 & 0 / 3 / 5 / 6 & 3 / 5 / 5 / 1 \\
    Average & 40 / 55 & 1 / 3 / 6 / 6 & 3 / 6 / 6 / 2 \\

\end{longtable}
}

The comparison indicates that AI scoring aligns most closely 
with the average of both human raters in the lower-to-mid 
score ranges (0–2), while slight divergences appear in the 
highest score category (Score~3). 
This observation is consistent with prior findings that LLMs 
exhibit mild tendencies toward under- or over-scoring nuanced, 
high-quality responses, often depending on the interpretive 
richness or conciseness of the answer.

\subsection{Quality of LLM Scoring Explanations}

A qualitative thematic analysis was performed on 200 randomly 
selected scoring explanations generated by the LLMs. 
Each explanation accompanied the numerical score and detailed 
the reasoning behind the model’s decision, including keyword 
matching, conceptual understanding, and alignment with rubric 
criteria. Each explanation was independently coded and rated 
according to three dimensions: \textit{clarity}, 
\textit{rubric alignment}, and \textit{coverage of evaluation criteria}.  
Results showed that 87\% of explanations were rated as high quality, 9\% 
as moderate, and only 4\% as low.  
These findings suggest that the LLMs not only achieved 
numerical accuracy but also provided coherent and 
interpretable justifications that aligned with human reasoning 
patterns.

\subsection{Frequency and Nature of LLM Scoring Errors}

Identified LLM scoring errors were classified into four 
primary categories to better understand performance 
limitations:

\begin{enumerate}
    \item \textbf{Under-scoring} — Correct answers scored too low (5.8\% of total cases).  
    \item \textbf{Over-scoring} — Incorrect answers scored too high (6.3\%).  
    \item \textbf{Partial-credit inconsistency} — Inconsistent scoring of partially correct responses (3.4\%).  
    \item \textbf{Rubric misinterpretation} — Misunderstanding or misapplication of rubric criteria (2.1\%).  
\end{enumerate}

Most errors were concentrated in questions belonging to the 
\textit{Analyzing} and \textit{Creating} categories, 
reaffirming that higher-order cognitive levels remain more 
difficult for LLMs to assess accurately. 
This emphasizes the importance of combining rubric-based 
prompting with rule-based validation or post-processing in 
future implementations to mitigate such inconsistencies. 

In summary, the evaluation metrics collectively demonstrate 
that the proposed framework achieves substantial agreement and 
interpretability in automated Thai-language grading. While 
performance approaches human-level reliability in lower and 
mid-level cognitive tasks, continued refinement in 
higher-order reasoning evaluation remains essential for 
pedagogical deployment.

